\documentclass[11pt]{article}

\usepackage{amsmath, amssymb, amsthm, graphicx, geometry, booktabs, hyperref, float, xcolor}

\geometry{margin=1in}

\title{Self-Organized Criticality in the Bak--Tang--Wiesenfeld Model}
\author{Kevin Wu \\ COMP90072 -- Complex Systems}
\date{\today}

\graphicspath{{figures/}{./}}

\newcommand{\Z}{\mathbb{Z}}
\newcommand{\N}{\mathbb{N}}
\newcommand{\calC}{\mathcal{C}}
\newcommand{\calS}{\mathcal{S}}

\newcommand\myworries[1]{\textcolor{red}{#1}}

% Safe figure command: compiles even if the image file is missing.
\newcommand{\plotfigure}[3]{
\begin{figure}[H]
    \centering
    \IfFileExists{#2}{
        \includegraphics[width=0.78\textwidth]{#2}
    }{
        \fbox{
        \parbox[c][5cm][c]{0.78\textwidth}{
            \centering
            Missing figure: \texttt{#2}\\[0.5em]
            #3
        }}
    }
    \caption{#1}
\end{figure}
}

\begin{document}

\maketitle

\begin{abstract}
This report studies self-organized criticality in the Bak--Tang--Wiesenfeld sandpile model. 
The model contains two-dimensional finite lattice and driven by adding grains of sand, and unstable sites topples grains to their nearest neighbours. 
The simulation records avalanche statistics including the number of topples, avalanche area, boundary loss, and avalanche length. 
The results are interpreted in terms of its power law like avalanche distributions. 
The report also discusses there results and the Abelian property of the model.
\end{abstract}

\tableofcontents
\newpage

% ----------------------
\section{Introduction}

\subsection{Background}

Complex systems are systems composed of many components that interact with each other. 
Examples includes Earth's global climate, organisms, and the human brain.

The Bak–Tang–Wiesenfeld (BTW) sandpile model is a simple example of such a system. 
The model consists of a lattice of sites, each containing a number of grains of sand. 
Grains are slowly added to the lattice. 
When the number of grains at a site reaches a threshold, that site topples and redistributes grains to neighbouring sites. 
This then can trigger a cascade of further topples called an avalanche. 
Although the microscopic update rule is simple, the model produces avalanches of many different sizes and therefore gives a clear setting in which to study self-organized criticality.

\subsection{Self-Organized Criticality}

Self-organized criticality (SOC) describes the tendency of certain kind of systems(typically slowly driven) to evolve toward a critical state without the need to finely tune parameters of the system. 

The key features of SOC are:

\begin{itemize}
    \item \textbf{Attractor states.} After an initial phase, the system approaches a statistically stable state. 
    The configuration continue to change, but statistics such as average height and avalanche distributions converges.

    \item \textbf{Scale invariance.} Events occur over many different scales. 
    There is no single characteristic avalanche size.

\end{itemize}

\subsection{Aim}

The aim of this project is to simulate the BTW model and study whether it exhibits the expected features of SOC. 
This includes:

\begin{itemize}
    \item to implement the standard BTW model on a finite two-dimensional lattice;
    \item to measure statistics such as topples, area, loss, and length;
    \item to compare the observed avalanche distributions with the expected behaviour;
    \item to discuss extensions of the model and predict how they affect the statistics.
\end{itemize}

\subsection{Solution to Question 1: Scale Invariance and Power Laws}

A scale-invariant distribution is one whose shape is unchanged by rescaling. 
Suppose \(P(x)\) is the frequency density of events of size \(x\). 
Scale invariance means that replacing \(x\) by \(\lambda x\) changes the distribution only by a multiplicative factor depending on \(\lambda\):
\[
    P(\lambda x) = \lambda^{\Delta}P(x).
\]
Setting \(x=1\) gives \(P(\lambda)=\lambda^{\Delta}P(1)\), so
\[
    P(\lambda x)=\frac{P(\lambda)P(x)}{P(1)}.
\]
Thus \(P\) is multiplicative up to normalization. The only continuous distribution that satisfy this condition is 
that it's logarithm is uniformally distributed.
Therefore a scale invariant system is associated with the power law.


\section{Bak--Tang--Wiesenfeld Model}

\subsection{Model Definition}

The simulation uses a finite rectangular lattice
\[
    \Gamma = \{1,\ldots,N\}\times \{1,\ldots,M\}.
\]
Each site \(x=(i,j)\in \Gamma\) has an integer height
\[
    p_t(x) = p_t(i,j) \in \Z_{\geq 0},
\]
which represents the number of grains at that site.

At each time step \(t\), one grain is added to a randomly selected site:
\[
    p_t(i,j) \leftarrow p_{t-1}(i,j)+1.
\]
If a site reaches the threshold
\[
    p_t(i,j) \geq 4,
\]
then it is unstable. 
A toppling at an unstable site removes four grains from that site and sends one grain to each of the four neighbours:
\[
    p_t(i,j) \leftarrow p_t(i,j)-4,
\]
\[
    p_t(i\pm1,j\pm1) \leftarrow p_t(i\pm1,j\pm1)+1,
\]
whenever these sites exist in the lattice.

If a grain is sent outside the lattice, it is removed from the system. 

\subsection{Solution to Question 2: Why the BTW Model Exhibits SOC}
The BTW model satisfies the main requirements for SOC in the following ways.

\begin{itemize}
    \item It's a slowly driven system as the external driving force is just the addition of a single grain.
    \item It self organizes as no external parameter is tuned to a critical value.
    \item It is dissipative as grains can leave the system at the boundary
\end{itemize}


\subsection{Solution to Question 3: Importance of Finite Boundaries}

Finite open boundaries are important because they provide dissipation. 
Every time a boundary site topples, some grains that would have gone to off-lattice neighbours are deleted.  This allows the system to balance input and output.

Another reason is that it takes too much resources to simulate infinite boundaries, as the sand will continue to spread out and takes up more space.
\subsection{Avalanches}

An avalanche is the chain reaction caused by adding one grain. 
The simulation records four avalanche features:

\begin{itemize}
    \item \textbf{Topples.} The total number of topplings in the avalanche.

    \item \textbf{Area.} The number of distinct sites that topple at least once.

    \item \textbf{Loss.} The total number of grains lost through the open boundary.

    \item \textbf{Length.} The spatial extent of the avalanche. 
    In this implementation, length is defined as the maximum Manhattan distance from the grain-addition site to any toppled site:
    \[
        \ell = \max_{x\in A} \|x-x_0\|_1
        =
        \max_{(i,j)\in A} \left(|i-i_0|+|j-j_0|\right),
    \]
    where \(x_0=(i_0,j_0)\) is the grain-addition site and \(A\) is the set of toppled sites.
\end{itemize}

The Manhattan metric is a natural choice because the BTW model only allows motion along horizontal and vertical nearest neighbour edges. 
It therefore matches the geometry of the toppling rule better than Euclidean distance.

\subsection{Solution to Question 4: Expected Average Height}

The average height of the lattice is
\[
    \langle p_t \rangle
    =
    \frac{1}{NM}\sum_{i, j}^{} p_t(i,j).
\]
It's clear that we have an upper bound of 3. Since each site has heights of 0, 1, 2, 3.
Assume evenly distributed, the average height would be
\[
    \frac{0 + 1 + 2 + 3}{4} = 1.5.
\]
However this assumption isn't true as adding sand will result the larger heights being more common in the system. So a more reasonable guess will be around $2$ (Turned out to be around 2.1).

\subsection{Steady State and Stationarity}

Starting from an empty lattice, the system first undergoes a transient phase. 
During this phase sand is dropped but not much toppling is occurring due to low average height of the system. Not much grains are lost at this phase.

Eventually, avalanches become frequent enough that the average boundary loss balances the one grain added per step. The system then enters a statistically stationary state. 

\subsection{Solution to Question 5: Steady State, Statistical Stationarity, and Equilibrium}

A \textbf{steady state} means that the dynamic condition where system properties remain constant over time due to balanced input and output.

A \textbf{statistically stationary state} means that statistical quantities of features do not change over time. 

An \textbf{equilibrium state} means there is no net flow through the system and that detailed balance holds. 
In an equilibrium system, there is no persistent input and output of mass or energy.

The BTW model is not an equilibrium system because of the continuous additions of grains and dissipation changing the layout.
It is better described as a statistically stationary state, in which the following quantities:

\begin{itemize}
    \item average height;
    \item average loss per grain addition;
    \item avalanche toppling distribution;
    \item avalanche area distribution;
    \item avalanche length distribution;
\end{itemize}

remains approximately constant.

\subsection{Solution to Question 6: Time Sampling, Ensemble Sampling, and Ergodicity}

Sampling a system over time means observing one long simulation trajectory.
Sampling an ensemble means running many independent systems and observing their configurations at comparable stages.

These two procedures are equivalent only under an ergodicity assumption. (Not Done)

\section{Abelian Property}

\subsection{Definition}

The BTW model is Abelian because the final stable configuration does not depend on the order in which unstable sites are toppled. 
If multiple sites are unstable at the same time, one may topple them in any order and still arrive at the same final stable state.

\subsection{Mathematical Formulation}

Let \(\Gamma\subset \Z^2\) be the finite lattice, and let
\[
    U_\Gamma = \{z:\Gamma\to \Z_{\geq 0}\}
\]
be the configuration space. 
The stable configurations form the subset
\[
    S_\Gamma =\{p\in\calC: p(x)<4 \text{ for all } x\in\Gamma\}.
\]

For \(x,y\in \Gamma\), define the BTW toppling matrix by
\[
    \Delta(x,y)
    =
    \begin{cases}
        -2d, & x=y,\\
        1, & x\text{ and }y\text{ are nearest neighbours},\\
        0, & \text{otherwise}.
    \end{cases}
\]
For a finite lattice with open boundaries, neighbours outside the lattice are not included. 
Thus boundary topplings lose grains.

A toppling at site \(x\) maps a configuration \(z\) to a new configuration \(T_x z\) given by
\[
    (T_xp)(y)=p(y)+\Delta(x,y),
\]

\subsection{Solution to Question 9: Meaning of the Toppling Matrix Conditions}

The general toppling matrix conditions are:

\[
    \Delta(x,y)=\Delta(y,x)\geq 0
    \quad \text{for } x\neq y,
\]
\[
    \Delta(x,x)<0,
\]
\[
    \sum_{y\in\Gamma}\Delta(x,y)\leq 0,
\]
\[
    \sum_{x,y\in\Gamma}\Delta(x,y)<0.
\]

Their meanings are as follows.

\begin{itemize}
    \item The first equation means two things. The  $\Delta(y,x)\geq 0$ part states that if the point is not the toppling point, the height should not decrease. The $\Delta(x,y)=\Delta(y,x)$ part states that if topping at point $x$ affects point $y$, toppling at point $y$ should affect point $x$ as well. These are all core properties of the model.

    \item  The second condition \(\Delta(x,x)<0\) means that toppling at \(x\) reduces the height at \(x\).  Without this, a toppling would not stabilize the unstable site.

    \item The condition
    \[
        \sum_y \Delta(x,y)\leq 0
    \]
    means that toppling does not create grains. 

    \item The condition
    \[
        \sum_{x,y}\Delta(x,y)<0
    \]
    means that somewhere in the system there is dissipation. 
    For the open-boundary BTW model, this dissipation occurs at the boundary. 
\end{itemize}

\subsection{Solution to Question 10: BTW Toppling Matrix and Threshold}

For the two-dimensional square lattice, each interior site has four nearest neighbours. 
The BTW toppling matrix is
\[
    \Delta(x,y)
    =
    \begin{cases}
        -2d, & x=y,\\
        1, & x\sim y,\\
        0, & \text{otherwise}.
    \end{cases}
\]
We are only considering the case for $d = 2$, however this generalizes to all $d > 2$.
This means that toppling at \(x\) removes $2d$ grains from \(x\) and adds one grain to each neighbour.

The stability threshold is
\[
    p_c = 2d.
\]
A site is stable if \(p(x)<2d\), and unstable if \(p(x)\geq 2d\).

The matrix satisfies the required conditions:

\begin{itemize}
    \item If \(x\neq y\), then \(\Delta(x,y)\) is either \(1\) for neighbours or \(0\) otherwise, so it is non-negative. 
    It is symmetric because the neighbouring property is symmetric.

    \item The diagonal entries are \(\Delta(x,x)=-2d<0\).

    \item For an interior site, the sum is
    \[
        -2d+1+1+1+1>0.
    \]
    For an edge site, the sum is negative because fewer than four neighbours lie inside the lattice. 
    For example, a non-corner edge site has sum
    \[
        -2d+3=-1.
    \]
    A corner has sum
    \[
        -2d+2=-2.
    \]
    Therefore \(\sum_y \Delta(x,y)\leq 0\).

    \item Since boundary sites have negative sum, the total sum over the finite lattice is strictly negative:
    \[
        \sum_{x,y}\Delta(x,y)<0.
    \]
\end{itemize}

\subsection{Solution to Question 11: Commutation of Toppling Operators}

Let \(T_x\) and \(T_y\) be toppling operators at sites \(x\) and \(y\). 
Assume both topplings are legal for the current unstable configuration \(z\). 
Then
\[
    (T_xz)(w)=z(w)+\Delta(x,w),
\]
and
\[
    (T_yT_xz)(w)
    =
    z(w)+\Delta(x,w)+\Delta(y,w).
\]
Similarly,
\[
    (T_xT_yz)(w)
    =
    z(w)+\Delta(y,w)+\Delta(x,w).
\]
Since addition is commutative,
\[
    z(w)+\Delta(x,w)+\Delta(y,w)
    =
    z(w)+\Delta(y,w)+\Delta(x,w).
\]
Therefore
\[
    T_yT_xz=T_xT_yz.
\]


\subsection{Solution to Question 12: Well-Defined Stabilization Operator}

\myworries{Need to add solution to this}


\section{Simulation Methodology}

\subsection{Implementation}

The simulation is implemented using an Python class for the Abelian sandpile. 
The lattice heights are stored in a two-dimensional NumPy array. 
The code is designed for large simulations such as \(100\times 100\) lattices over approximately \(100000\) grain additions.

The main implementation choices are:

\begin{itemize}
    \item \textbf{Grid storage.} The sandpile is stored as an integer array \texttt{grid} of shape \((W,H)\).

    \item \textbf{Threshold.} The default threshold is \(4\), matching the standard two-dimensional BTW model.

    \item \textbf{Neighbour precomputation.} The valid neighbours of each lattice site are precomputed, avoiding repeated boundary checks.

    \item \textbf{Boundary loss precomputation.} For each site, the number of missing neighbours is precomputed. 
    This determines how many grains are lost when that site topples.

    \item \textbf{Unstable-site set.} The code maintains a set of unstable sites rather than repeatedly scanning the whole grid.

    \item \textbf{Batch toppling.} If a site has height \(p(i,j)\), it is toppled
    \[
        q=\left\lfloor \frac{p(i,j)}{4}\right\rfloor
    \]
    times in one operation. 
    This is faster than toppling the same site one step at a time.
\end{itemize}

\subsection{Algorithm}

For each simulation step:

\begin{enumerate}
    \item Select a site \(x_0=(i_0,j_0)\), usually uniformly at random.
    \item Add one grain to \(x_0\).
    \item If \(x_0\) is unstable, add it to the unstable set.
    \item While the unstable set is non-empty:
    \begin{enumerate}
        \item remove an unstable site \(x\);
        \item compute \(q=\lfloor z(x)/4\rfloor\);
        \item subtract \(4q\) grains from \(x\);
        \item add \(q\) grains to each valid neighbour;
        \item add the corresponding boundary loss;
        \item record that \(x\) toppled;
        \item add any newly unstable neighbours back into the unstable set.
    \end{enumerate}
    \item Once all sites are stable, record the avalanche statistics.
\end{enumerate}

The Abelian property guarantees that the final stable configuration is independent of the order in which the unstable set is processed. 

\subsection{Data Collection}

For each avalanche, the code records:

\begin{itemize}
    \item \textbf{Topples:} increased by \(q\) whenever a site batch-topples \(q\) times.

    \item \textbf{Area:} computed as the number of distinct sites that toppled at least once.

    \item \textbf{Loss:} increased by \(q\) times the number of missing neighbours at the toppled site.

    \item \textbf{Length:} computed as the maximum Manhattan distance from the grain-addition site to any toppled site.
\end{itemize}

These values are stored as arrays or a table so that they can be plotted after the simulation. 

\section{Results}

\subsection{Avalanche Distributions}

\plotfigure
{Distribution of avalanche topples.}
{topples_log_binned_density.png}

\plotfigure
{Distribution of avalanche area.}
{area_log_binned_density.png}

\plotfigure
{Distribution of boundary loss.}
{loss_log_binned_density.png}

\plotfigure
{Distribution of avalanche length.}
{length_log_binned_density.png}

The avalanche distributions are strongly left-skewed.  Most grain added induces small avalanches and occasionally induces large ones. The distribution is consistent with the power law:
\[
    P(x)\sim x^{-\alpha}.
\]
In a finite lattice, the power-law behaviour cannot continue indefinitely. 
Due to the limiting grid size, the tail of the power function has been cut off(As evident in the plots).
The key evidence of the power law is the linear segment in the middle of the log log plot ignoring the cutoff at the very end.

\subsection{Solution to Question 7: Distribution Type and Fitting}

The recorded avalanche frequencies should be plotted for topples, area, loss, and length. 
The expected results are:

\begin{itemize}
    \item \textbf{Topples:} heavy-tailed and approximately power-law-like, with many small avalanches and rare large avalanches.

    \item \textbf{Area:} also heavy-tailed and power-law-like, though with a finite-size cutoff at \(WH\).

    \item \textbf{Loss:} often has many small or zero values because avalanches must reach the boundary to lose grains. 
    Its distribution is still typically right-skewed.

    \item \textbf{Length:} heavy-tailed up to the maximum possible lattice distance.
\end{itemize}

A simple power-law fit can be obtained by fitting a line to the log log frequency plot:
\[
    \log P(x)=c-\alpha \log x.
\]
The exponent is then the negative slope:
\[
    \alpha=-\text{slope}.
\]
\myworries{Need to find fit}


\subsection{Scaling Behaviour}

The results support scale invariance if avalanches occur over many sizes and the log--log plots show approximately linear scaling regions. 

We can also construct plots to compare the results from simulation from different sizes.

\myworries{Add graphs comparing the stats from different sizes plotted onto 1 graph}


\subsection{Correlations}

\plotfigure
{Scatter plot of avalanche area against number of topples.}
{area_vs_topples.png}

\plotfigure
{Scatter plot of avalanche length against avalanche area.}
{length_vs_area.png}

\subsection{Solution to Question 8: Expected and Observed Correlations}

The four avalanche features are expected to be correlated, but not identical.

\begin{itemize}
    \item \textbf{Topples and area} should be strongly positively correlated. 
    A larger avalanche usually involves more sites and more topplings. 

    \item \textbf{Area and length} should also be positively correlated. 
    Avalanches that spread farther tend to visit more sites. 

    \item \textbf{Loss and length} should be positively correlated because grains are lost only when avalanches reach the boundary. 
    Avalanches with higher length are more likely to interact with the boundary.

    \item \textbf{Loss and topples} should also be positively correlated, since larger avalanches have more opportunities to reach the boundary and lose grains.
\end{itemize}



\section{Model Extensions}

\section{Conclusion}

I'll have to think of one


\begin{thebibliography}{9}

\bibitem{btw1987}
P. Bak, C. Tang, and K. Wiesenfeld,
\emph{Self-organized criticality: An explanation of the 1/f noise},
Physical Review Letters, 59(4), 381--384, 1987.

\bibitem{btw1988}
P. Bak, C. Tang, and K. Wiesenfeld,
\emph{Self-organized criticality},
Physical Review A, 38(1), 364--374, 1988.

\bibitem{dhar1990}
D. Dhar,
\emph{Self-organized critical state of sandpile automaton models},
Physical Review Letters, 64(14), 1613--1616, 1990.

\bibitem{newman2005}
M. E. J. Newman,
\emph{Power laws, Pareto distributions and Zipf's law},
Contemporary Physics, 46(5), 323--351, 2005.

\bibitem{soc_sciencedirect}
Elsevier,
\emph{Self-Organized Criticality},
ScienceDirect Topics.
Available at:
\url{https://www.sciencedirect.com/topics/materials-science/self-organized-criticality},
accessed 15 May 2026.

\bibitem{powerlaw_scale}
Math StackExchange,
\emph{Scale invariance of power law distribution}.
Available at:
\url{https://math.stackexchange.com/questions/140962/scale-invariance-of-power-law-distribution},
accessed 15 May 2026.

\bibitem{bak_review}
P. Bak, K. Christensen, L. Danon, and T. Scanlon,
\emph{Unified scaling law for earthquakes},
arXiv preprint cond-mat/0207674, 2002.
Available at:
\url{https://arxiv.org/pdf/cond-mat/0207674}.

\bibitem{todo_notes}
TeX StackExchange,
\emph{How to add TODO notes?}
Available at:
\url{https://tex.stackexchange.com/questions/9796/how-to-add-todo-notes},
accessed 15 May 2026.

\end{thebibliography}

\end{document}
